% titel: Blatt 10, A2: Zunehmende Säurestärke ordnen
% Teil II: der durchdachte Weg.

\subsection*{a) HF, HI, NaH, H$_2$O}
\begin{itemize}
  \item \textbf{NaH} ist kein Brønsted-Protonendonator (liefert H$^-$), also
        die schwächste ``Säure`` --- praktisch keine Säure, starke Base.
  \item \textbf{H$_2$O} gibt nur sehr schwer ein Proton ab (schwache Säure).
  \item \textbf{HF vs.\ HI:} die H--X-Bindung wird mit zunehmender Atomgrösse
        schwächer/stabiler ginge es also umgekehrt --- entscheidend ist die
        Stabilität der Base. I$^-$ (gross, polarisierbar) ist die stabilste
        Base, F$^-$ die reaktivste. Also ist \textbf{HI die stärkere,
        HF die schwächere} der beiden.
\end{itemize}
Reihenfolge (zunehmend): NaH $<$ H$_2$O $<$ HF $<$ HI.

\subsection*{b) H$_2$SeO$_3$, H$_2$SeO$_4$, H$_2$O}
Unter den Sauerstoffsäuren macht mehr O pro Zentralatom die Säure stärker
(stabileres Anion). H$_2$SeO$_4$ (mit mehr O) ist die stärkste.
Reihenfolge: H$_2$O $<$ H$_2$SeO$_3$ $<$ H$_2$SeO$_4$.

\begin{ergebnis}
a) NaH $<$ H$_2$O $<$ HF $<$ HI;
b) H$_2$O $<$ H$_2$SeO$_3$ $<$ H$_2$SeO$_4$.
So steigt die Säurestärke.
\end{ergebnis}

\begin{kontrolle}
Erwartungswerte passen zur Theorie: unter den Halogenwasserstoffen steigt
die Stärke mit der Grösse der Halogene; mehr O im Anion (SeO$_4^{2-}$ vs.
SeO$_3^{2-}$) verstärkt die Säure. NaH als reine Base steht logischerweise
am Anfang.
\end{kontrolle}

\begin{merksatz}
Säurestärke: bei HX wächst sie mit der Stabilität des Anions (HF am
schwächsten), bei O-Säuren mit der Zahl der O-Atome am Zentralatom.
Metallhydride sind keine Säuren.
\end{merksatz}
